% ============================================================
%  CHAPTER 7 — Manipulation and Grasping Pipeline
% ============================================================

\chapter{Manipulation and Grasping Pipeline}
\label{ch:manipulation}

\section{Overview}

Grasping an object from a table is the core physical capability of the system.
Despite appearing simple, reliable grasping on a real robot requires solving several
tightly coupled sub-problems simultaneously: detecting the object and estimating its
3-D pose, reasoning about the table surface, planning a collision-free arm trajectory,
and executing the motion with sufficient precision for the gripper to close around the
object.

This chapter describes the full grasping pipeline developed for TIAGo, from raw RGB-D
sensor data to a closed gripper holding the object.
The pipeline is implemented in \texttt{reach\_object\_v5\_torso\_descent\_working.py}
and is invoked by the \texttt{GrabBottleSkill} inside the agent.
Although the system was initially developed to grasp bottles, the detection and motion
pipeline is fully object-agnostic: any object detectable by YOLO11n can be targeted
by changing the query class.

\section{Pipeline Evolution}

The grasping pipeline was developed iteratively across six versions.
Each version addressed a specific failure mode identified during real-robot testing.
Table~\ref{tab:grasp_versions} summarises the evolution.

\begin{table}[H]
\centering
\caption{Evolution of the grasping pipeline across versions.}
\label{tab:grasp_versions}
\begin{tabularx}{\linewidth}{lXX}
\toprule
\textbf{Version} & \textbf{Key change} & \textbf{Problem addressed} \\
\midrule
v1 & Basic MoveIt goal, YOLOv8, \texttt{moveit\_commander} &
     First working prototype; no depth, no pose estimation \\
v2 & 5-frame detection averaging, aligned depth &
     Noisy single-frame detections caused arm to miss object \\
v3 & PAL gripper integration, calibrated offsets &
     Gripper closed in wrong position relative to arm\_tool\_link \\
v4 & RANSAC cylinder fit on depth point cloud &
     Bounding-box centroid inaccurate for thin objects; cylinder gives true centre \\
v5 & Torso descent for final approach, YOLO11n via HTTP service &
     MoveIt Cartesian path failed at low wrist heights; torso descent avoids singularities \\
v6 & MoveJ + MoveL with LSPB velocity profile (planned) &
     Torso descent is slow and produces an arc path; replace with proper trajectory generation \\
\bottomrule
\end{tabularx}
\end{table}

\section{Full Pipeline: Version 5}

The complete pipeline executed by \texttt{BottleReacher.run()} consists of nine
sequential steps, illustrated in Figure~\ref{fig:grasp_pipeline}.

\begin{figure}[H]
\centering
\begin{tikzpicture}[
  step/.style={rectangle, rounded corners=3pt, draw=darkblue, thick,
               fill=lightblue!30, minimum width=8cm, minimum height=0.7cm,
               align=left, font=\small, text width=7.5cm},
  arr/.style={->, thick, draw=darkblue},
  lbl/.style={font=\small\bfseries\color{darkblue}, minimum width=1.0cm, align=right},
  node distance=0.35cm
]
  \node[pstep] (s1) {\textbf{Step 1} \quad Tilt head to $-0.5$ rad (looking at table), preserve pan};
  \node[pstep, below=of s1] (s2) {\textbf{Step 2} \quad Unfold arm (\texttt{play\_motion('unfold\_arm')})};
  \node[pstep, below=of s2] (s3) {\textbf{Step 3} \quad Open gripper (\texttt{play\_motion('open')})};
  \node[pstep, below=of s3] (s4) {\textbf{Step 4} \quad Get camera TF once (tf\_echo subprocess)};
  \node[pstep, below=of s4] (s5) {\textbf{Step 5} \quad Detect object — YOLO11n + RANSAC cylinder fit, 5-frame average};
  \node[pstep, below=of s5] (s6) {\textbf{Step 6} \quad Add table and object collision boxes to MoveIt scene};
  \node[pstep, below=of s6] (s7) {\textbf{Step 7} \quad OMPL plan to pre-grasp (above object)};
  \node[pstep, below=of s7] (s8) {\textbf{Step 8} \quad Torso descent (vertical straight line to grasp)};
  \node[pstep, below=of s8] (s9) {\textbf{Step 9} \quad Close gripper + raise torso (lift object)};

  \draw[arr] (s1) -- (s2);
  \draw[arr] (s2) -- (s3);
  \draw[arr] (s3) -- (s4);
  \draw[arr] (s4) -- (s5);
  \draw[arr] (s5) -- (s6);
  \draw[arr] (s6) -- (s7);
  \draw[arr] (s7) -- (s8);
  \draw[arr] (s8) -- (s9);
\end{tikzpicture}
\caption{Sequential steps of the v5 grasping pipeline.}
\label{fig:grasp_pipeline}
\end{figure}

\section{Object Detection and 3-D Pose Estimation}
\label{sec:grasp_perception}

\subsection{YOLO11n Detection}

At each perception frame, the current RGB image is encoded as JPEG and sent to the
YOLO11n HTTP microservice (port 5001) with an \texttt{X-Classes} header specifying
the target class.
The service returns bounding boxes $[x, y, w, h]$ in pixel coordinates with associated
confidence scores.
The detection with the highest confidence above a threshold $\tau = 0.4$ is selected.

\subsection{Depth Point Cloud Extraction}

Rather than using only the bounding box centroid, the system projects all depth pixels
within the YOLO bounding box into 3-D space.
A 4-pixel margin is applied to each edge of the box to discard boundary pixels, which
typically contain mixed foreground/background depth values.

For each valid pixel $(u, v)$ with depth $d$ (filtered to $0.15 < d < 3.5$ m), the
camera-frame 3-D point is computed using the pinhole model:

\begin{equation}
  \mathbf{p}_{\text{cam}} =
  \begin{pmatrix}
    \dfrac{(u - c_x) \cdot d}{f_x} \\[8pt]
    \dfrac{(v - c_y) \cdot d}{f_y} \\[8pt]
    d
  \end{pmatrix}
\end{equation}

where $(f_x, f_y, c_x, c_y)$ are the camera intrinsics from \texttt{/xtion/rgb/camera\_info}.
Each point is then transformed into the \texttt{base\_footprint} frame:

\begin{equation}
  \mathbf{p}_{\text{base}} = \mathbf{R}_{\text{cam}}^{\text{base}} \, \mathbf{p}_{\text{cam}}
                            + \mathbf{t}_{\text{cam}}^{\text{base}}
\end{equation}

where $\mathbf{R}$ and $\mathbf{t}$ are retrieved once at the start of each grasp
attempt via a \texttt{rosrun tf tf\_echo} subprocess call.

\subsection{RANSAC Cylinder Fitting}
\label{sec:ransac}

The projected 3-D point cloud is fitted with a vertical cylinder (axis parallel to
world $Z$) using RANSAC (Random Sample Consensus).
The cylinder model is defined by its centre $(c_x, c_y)$ and radius $r$ in the $XY$
plane; the axis is constrained to be vertical so no axis direction needs to be estimated.

\textbf{RANSAC iteration:}
\begin{enumerate}[noitemsep]
  \item Sample 3 random points from the $XY$ projection of the cloud.
  \item Fit a circle through the 3 points using the circumscribed circle formula:
        \begin{align}
          D &= 2\bigl[a_x(b_y - c_y) + b_x(c_y - a_y) + c_x(a_y - b_y)\bigr] \notag \\
          u_x &= \frac{(a_x^2+a_y^2)(b_y-c_y) + (b_x^2+b_y^2)(c_y-a_y) + (c_x^2+c_y^2)(a_y-b_y)}{D}
        \end{align}
  \item Reject if $r \notin [0.015, 0.060]$ m (not a plausible object radius).
  \item Count inliers: points within 15 mm of the circle circumference.
  \item Keep the hypothesis with the most inliers.
\end{enumerate}

After 300 iterations, the best hypothesis is refined using algebraic least-squares on
all inliers.
The fit is accepted only if the inlier count exceeds 8 and the inlier ratio exceeds
$30\%$.
The estimated quantities are:

\begin{center}
\begin{tabular}{ll}
\toprule
\textbf{Output} & \textbf{Description} \\
\midrule
$c_x, c_y$ & Object centre in \texttt{base\_footprint} XY plane \\
$r$ & Estimated object radius (m) \\
$z_{\text{top}}$ & 95th percentile of inlier $Z$ values (object top) \\
$z_{\text{bottom}}$ & 5th percentile of inlier $Z$ values (object base) \\
$z_{\text{table}}$ & Table surface height (sampled below YOLO box) \\
\bottomrule
\end{tabular}
\end{center}

\subsection{Table Height Estimation}

The table surface height $z_{\text{table}}$ is estimated by sampling the depth image
in a horizontal strip of 8 pixels immediately below the bottom edge of the YOLO
bounding box — the region where the object meets the table surface.
The median depth of valid pixels in this strip is back-projected to obtain the table $Z$
coordinate in \texttt{base\_footprint}.
This estimate is used to position the MoveIt collision table at the correct height.

\subsection{5-Frame Averaging for Stability}

A single cylinder fit on one frame is sensitive to sensor noise and small camera
vibrations.
The system collects 5 successful fits ($\text{inlier\_ratio} > 0.3$) and averages the
key quantities:

\begin{equation}
  \hat{c}_x = \frac{1}{5}\sum_{i=1}^{5} c_x^{(i)}, \quad
  \hat{c}_y = \frac{1}{5}\sum_{i=1}^{5} c_y^{(i)}, \quad
  \hat{z}_{\text{top}} = \frac{1}{5}\sum_{i=1}^{5} z_{\text{top}}^{(i)}
\end{equation}

The standard deviation across the 5 samples is logged as a quality metric.
Typical values for a stationary object are $\sigma_{c_x} < 3$ mm, indicating
stable pose estimation.

\section{MoveIt Planning Scene}

Before planning arm motion, two collision objects are added to the MoveIt planning scene:

\textbf{Table:} A flat box of dimensions $0.8 \times 1.2 \times 0.02$ m, centred at
the estimated table surface position.
This prevents the arm from planning through the table surface.

\textbf{Object:} A vertical cylinder with the estimated radius plus a 1 cm safety
margin and height $z_{\text{top}} - z_{\text{bottom}}$.
This is added before the OMPL plan to the pre-grasp position so that the planner
routes around the object.
It is \emph{removed} before the descent phase because the gripper must intentionally
enter the object's bounding volume to close around it.

\begin{lstlisting}[language=Python, caption={Collision object lifecycle during grasp.}]
# Before OMPL plan:
self.add_table_collision(bx, z_table - 0.01)
self.add_object_collision(cx, cy, z_bottom, z_top, radius)

# OMPL plans around both → arm reaches above object
self._send_arm_goal(cx, cy, z_top + 0.20)   # pre-grasp 20 cm above

# Before descent:
self.remove_collision_object("object")       # allow gripper to enter

# Torso descends straight down → gripper at grasp height
# Gripper closes → object grasped
\end{lstlisting}

\section{Arm Motion: OMPL + Torso Descent}

\subsection{OMPL Plan to Pre-Grasp}

The OMPL (Open Motion Planning Library) planner, accessed via the MoveIt
\texttt{/move\_group} action server, plans a collision-free trajectory from the current
arm configuration to a pre-grasp pose 20 cm directly above the estimated object centre.
The goal is expressed as a Cartesian position constraint (5 cm sphere) on the
\texttt{arm\_tool\_link} with a soft orientation constraint ($\pm 0.2$ rad tolerance)
to maintain the grasp orientation.

Key planning parameters:
\begin{center}
\begin{tabular}{ll}
\toprule
\textbf{Parameter} & \textbf{Value} \\
\midrule
Planning group & \texttt{arm\_torso} \\
Planning attempts & 10 \\
Allowed planning time & 10 s \\
Replan attempts & 3 \\
Max velocity scaling & 0.3 \\
Max acceleration scaling & 0.2 \\
\bottomrule
\end{tabular}
\end{center}

\subsection{Torso Descent — Vertical Approach}
\label{sec:torso_descent}

Once the arm is positioned at the pre-grasp height, the final 20 cm vertical
approach is performed by lowering the torso lift joint rather than re-planning with
MoveIt.

This design choice is motivated by a key property of TIAGo's kinematic structure:
the torso lift joint is a \textbf{prismatic joint with its axis aligned with world $Z$}.
Lowering the torso therefore moves the entire arm assembly — and consequently the
\texttt{arm\_tool\_link} — \emph{exactly vertically downward}, with no lateral deviation.
This produces a straight-line end-effector path without requiring Cartesian path
computation or IK solutions at each intermediate height.

The torso is lowered using the \texttt{/torso\_controller/follow\_joint\_trajectory}
action server with a trapezoidal velocity profile (single-point trajectory, 3 s duration):

\begin{lstlisting}[language=Python, caption={Torso descent implementation.}]
def _torso_descent(self, delta_z):
    target = np.clip(self.current_torso_height - delta_z, 0.0, 0.35)
    traj = JointTrajectory()
    traj.joint_names = ['torso_lift_joint']
    pt = JointTrajectoryPoint()
    pt.positions    = [target]
    pt.time_from_start = rospy.Duration(3.0)
    traj.points.append(pt)
    goal = FollowJointTrajectoryGoal()
    goal.trajectory = traj
    self.torso_client.send_goal(goal)
    self.torso_client.wait_for_result(rospy.Duration(10))
\end{lstlisting}

The current torso height is tracked continuously from \texttt{/joint\_states} so
the actual joint position is used rather than an assumed value.

\subsection{Gripper Close and Object Lift}

After the descent, \texttt{play\_motion('close')} closes the gripper.
The arm is then returned to a safe height by raising the torso back to 10 cm, lifting
the object clear of the table before any further motion is planned.

\section{Grasp Offsets and Calibration}

The relationship between the cylinder centre $(c_x, c_y, z_{\text{top}})$ estimated
from depth and the actual target position of \texttt{arm\_tool\_link} for a successful
grasp is determined by calibrated offsets $(\Delta x, \Delta y, \Delta z)$:

\begin{equation}
  \mathbf{p}_{\text{grasp}} =
  \begin{pmatrix} c_x + \Delta x \\ c_y + \Delta y \\ z_{\text{top}} + \Delta z \end{pmatrix}
\end{equation}

These offsets account for the physical offset between the \texttt{arm\_tool\_link}
origin and the gripper jaw centre, and are calibrated using
\texttt{calibrate\_grasp\_interactive.py} — an interactive tool that lets the user
adjust offsets while the robot holds a known position, then saves them to
\texttt{grasp\_offsets.yaml}.
Default values ($\Delta x = 0.16$, $\Delta y = -0.04$, $\Delta z = -0.04$ m) are used
if no calibration file exists.
Environment variables \texttt{GRASP\_DX}, \texttt{GRASP\_DY}, \texttt{GRASP\_DZ}
allow runtime override without editing files.

\section{Grasp Orientation}

The gripper orientation throughout the grasp is fixed at:

\begin{equation}
  \mathbf{q}_{\text{grasp}} = (x=0.703,\; y=0.073,\; z=-0.034,\; w=0.706)
\end{equation}

which corresponds to a \textbf{top-down approach} with the gripper fingers aligned
with the robot's forward direction.
This orientation is maintained as a soft constraint ($\pm 0.2$ rad tolerance on all
axes) during OMPL planning and held fixed during the torso descent.
A top-down approach is used because it generalises to any object placed on a flat
surface — the fingers straddle the object regardless of its shape.

\section{Head Control During Grasping}

During the grasping pipeline, the head is tilted downward to $-0.5$ rad (looking at
the table) to keep the object within the camera field of view.
The current pan angle is preserved by reading \texttt{head\_1\_joint} from
\texttt{/joint\_states} before sending the trajectory command:

\begin{lstlisting}[language=Python, caption={Head tilt with pan preservation.}]
def _tilt_head(self, tilt=-0.5, duration=2.0):
    js = rospy.wait_for_message('/joint_states', JointState, timeout=2.0)
    pan = js.position[js.name.index('head_1_joint')] \
          if 'head_1_joint' in js.name else 0.0
    # send trajectory with [pan, tilt] — pan unchanged, only tilt set
\end{lstlisting}

This ensures the head does not snap back to the left-right centre when only the
tilt needs to change — important when the object was found to one side by
\texttt{search\_with\_head}.

\section{Planned Improvement: Industrial-Style Trajectory (v6)}
\label{sec:v6_plan}

The torso descent approach (Section~\ref{sec:torso_descent}), while effective, has
two drawbacks:
\begin{itemize}[noitemsep]
  \item It is \textbf{slow}: lowering 20 cm at a safe speed takes approximately 3--4 s.
  \item The path is \textbf{kinematically coupled}: while the torso translates the arm
        endpoint vertically, the arm joints are locked, which can lead to subtle
        orientation drift if the arm is not perfectly positioned.
\end{itemize}

The planned v6 implementation adopts an industrial-style \textbf{MoveJ / MoveL} model:

\begin{description}[noitemsep]
  \item[\texttt{MoveJ}] Joint-space interpolation with a trapezoidal velocity profile
        for large repositioning moves (home $\rightarrow$ pre-grasp, lifted
        $\rightarrow$ home). The planner is free to choose the joint-space path.
  \item[\texttt{MoveL}] Cartesian straight-line motion for the precision approach and
        retreat, with a \textbf{Linear Segment with Parabolic Blends (LSPB)} velocity
        profile ensuring smooth acceleration and deceleration.
\end{description}

The LSPB profile for a segment of length $d$ with peak velocity $v_{\max}$ and
acceleration $a$:

\begin{equation}
  t_{\text{blend}} = \frac{v_{\max}}{a}, \quad
  t_{\text{total}} = \frac{d}{v_{\max}} + \frac{v_{\max}}{a}
\end{equation}

\[
  v(t) =
  \begin{cases}
    a \cdot t & 0 \leq t \leq t_{\text{blend}} \\
    v_{\max}  & t_{\text{blend}} < t \leq t_{\text{total}} - t_{\text{blend}} \\
    v_{\max} - a(t - (t_{\text{total}} - t_{\text{blend}})) & \text{otherwise}
  \end{cases}
\]

At via-points between two \texttt{MoveL} segments (e.g.\ the transition from vertical
descent to horizontal retreat), a \textbf{blending zone} is applied: the robot begins
transitioning to the next segment direction before reaching the exact via-point,
carving a smooth rounded corner rather than coming to a full stop.
This reduces cycle time and eliminates the velocity discontinuity at segment
boundaries.

The resulting v6 grasp motion sequence is:

\begin{center}
\begin{tabular}{lll}
\toprule
\textbf{Phase} & \textbf{Motion type} & \textbf{Speed} \\
\midrule
Home $\rightarrow$ pre-grasp  & MoveJ & 30\% of max \\
Pre-grasp $\rightarrow$ grasp & MoveL (straight down) & 3 cm/s \\
Gripper close                 & — & — \\
Grasp $\rightarrow$ lifted    & MoveL (straight up) & 5 cm/s \\
Lifted $\rightarrow$ retreat  & MoveL (blended with lift) & 8 cm/s \\
Retreat $\rightarrow$ home    & MoveJ & 30\% of max \\
\bottomrule
\end{tabular}
\end{center}

Since the motion is fully described by Cartesian poses, it is completely
\textbf{object-agnostic}: only the target $(c_x, c_y, z_{\text{top}})$ from the
perception pipeline changes between objects.
Orientation is held fixed throughout all \texttt{MoveL} phases at
$\mathbf{q}_{\text{grasp}}$ (no SLERP needed for a fixed top-down grasp).

\section{Summary}

Table~\ref{tab:grasp_summary} summarises the design decisions in the final v5
grasping pipeline.

\begin{table}[H]
\centering
\caption{Summary of key design decisions in the v5 grasping pipeline.}
\label{tab:grasp_summary}
\begin{tabularx}{\linewidth}{lXX}
\toprule
\textbf{Sub-problem} & \textbf{Solution} & \textbf{Rationale} \\
\midrule
Object detection & YOLO11n via HTTP service & Python 3.6 in Docker cannot run Ultralytics directly \\
3-D localisation & RANSAC cylinder fit on depth cloud & More accurate than bounding-box centroid for cylindrical objects \\
Pose stability   & 5-frame average & Reduces sensor noise; $\sigma < 3$ mm typical \\
Table height     & Depth strip below YOLO box & Direct measurement, more accurate than inference \\
Collision scene  & Table box + object cylinder & Prevents arm from colliding with surface \\
Approach path    & OMPL (free) + torso descent (vertical) & Avoids Cartesian singularities at low wrist heights \\
Descent path     & Torso prismatic joint & Guaranteed straight vertical line, no IK required \\
Orientation      & Fixed top-down quaternion & Generalises to any flat-surface object \\
Offsets          & Calibrated YAML + env var override & Accounts for tool-link to jaw-centre offset \\
\bottomrule
\end{tabularx}
\end{table}
